% =============================================================
% Chapitre 5 — Tests, validation, résultats et écarts (~10 pages)
% Le plan de tests (18 tests) est repris de la méthodologie projet
% (document réel) ; les RÉSULTATS restent [À COMPLÉTER] tant que
% 09_AUDIT_E2E_STAGING n'est pas fourni en intégralité — ne jamais
% remplir la colonne "résultat obtenu" par extrapolation.
% =============================================================
\chapter{Tests, validation, résultats et écarts}
\label{ch:tests}

\section{Méthodologie et plan de tests Zero Trust}

Le principe de validation retenu est celui du \emph{contrôle contourné volontairement} :
chaque test consiste à tenter de déjouer un contrôle de sécurité annoncé, et c'est l'échec
de cette tentative qui constitue la preuve. Dix-huit scénarios ont été définis, couvrant les
sept couches de l'architecture ainsi que les deux pipelines de livraison (applicatif et
infrastructure).

\begin{longtable}{@{}P{1cm}P{6.2cm}P{6.8cm}@{}}
\caption{Plan de tests de contournement Zero Trust}
\label{tab:plan-tests} \\
\toprule
\textbf{\#} & \textbf{Tentative} & \textbf{Résultat attendu} \\
\midrule
\endfirsthead
\toprule
\textbf{\#} & \textbf{Tentative} & \textbf{Résultat attendu} \\
\midrule
\endhead
T1 & Contournement du point d'entrée (accès direct à l'URL interne d'un service) & Refus
(403/404) \\
T2 & Connexion à la base de données applicative depuis Internet & Délai d'attente / refus ---
pas d'adresse publique \\
T3 & Appel du composant d'enrichissement sémantique sans jeton d'identité, depuis le réseau
privé & Refus (401/403) \\
T4 & Écriture dans la table de détections avec le compte de service de l'enrichissement &
Refus (intégrité des preuves) \\
T5 & Lecture d'un secret applicatif avec un compte de service d'un autre composant & Refus
(cloisonnement des secrets) \\
T6 & Création d'une clé d'accès statique pour un compte de service & Refus (politique
d'organisation) \\
T7 & Commit contenant un secret factice & Pipeline en échec à l'étape de détection de
secrets \\
T8 & Injection volontaire de code non filtré dans une proposition de modification & Pipeline
en échec à l'étape d'analyse statique \\
T9 & Image de base volontairement obsolète & Pipeline en échec à l'étape d'analyse des
vulnérabilités \\
T10 & Déploiement référencé par étiquette mutable plutôt que par empreinte & Refusé par
convention/contrôle du pipeline \\
T11 & Requête avec charge utile d'injection sur le point d'entrée & Refus par le pare-feu
applicatif, trace journalisée \\
T12 & Sortie réseau non déclarée depuis un service applicatif & Bloquée, journalisée \\
T13 & Scénario d'attaque simulé couvrant plusieurs tactiques sur une même entité en moins
d'une heure & Incident chaîné, preuves listées \\
T14 & Événement sans règle statique correspondante mais sémantiquement proche d'une
technique connue & Rattachement sémantique au-dessus du seuil, ou non-rattachement assumé \\
T15 & Modification manuelle hors du code d'infrastructure, puis nouvelle planification &
Dérive détectée et annoncée \\
T16 & Proposition d'infrastructure contenant une ressource volontairement mal configurée &
Pipeline infra en échec à l'étape d'analyse de sécurité \\
T17 & Tentative de déploiement vers la production depuis le pipeline de développement &
Refus (identité fédérée limitée au projet cible) \\
T18 & Comparaison des empreintes d'image entre les environnements & Identiques à l'octet
près \\
\bottomrule
\end{longtable}

\section{Résultats des tests de contournement}

\begin{attentedoc}
Cette section nécessite \texttt{09\_AUDIT\_E2E\_STAGING\_2026-08-07.md} en intégralité (ou
le rejeu documenté des 18 tests sur l'environnement de référence). Ne pas remplir la colonne
\og résultat obtenu \fg{} avant d'avoir cette source --- reprendre alors le même format à
deux temps (scénario/critère puis résultat/preuve) que celui utilisé au tableau~
\ref{tab:plan-tests}, en ajoutant une colonne \og preuve \fg{} (capture, extrait de
journal).
\end{attentedoc}

\section{Mesures avant/après}

Cinq indicateurs sont prévus pour comparer la posture de sécurité avant et après déploiement
du socle, sur un même jeu de journaux rejoué : le rappel de détection sur des scénarios
injectés, le taux de consolidation (alertes rapportées à des incidents), le délai moyen de
détection, le délai moyen de correction, et la couverture des techniques d'attaque
observables.

\begin{attentedoc}
Tableau à remplir avec des \textbf{mesures réelles} une fois disponibles --- ne jamais
présenter une estimation comme une mesure. Si une mesure n'est pas disponible au moment de
la rédaction, l'indiquer explicitement plutôt que l'omettre silencieusement.
\end{attentedoc}

\section{Écarts entre conception et implémentation réelle}

\begin{attentedoc}
Section à construire à partir de \texttt{06\_ECARTS\_IMPLEMENTATION.md}. Ne conserver du
chapitre~\ref{ch:mise-en-oeuvre} que ce qui est confirmé réellement déployé ; lister ici
chaque écart identifié avec sa raison.
\end{attentedoc}

\section{Difficultés rencontrées et corrections apportées}

Trois incidents réels, déjà consignés dans les documents projet, illustrent bien l'esprit du
guide pédagogique --- présenter les difficultés rencontrées plutôt qu'un déroulement
idéalisé :

\begin{itemize}
  \item \textbf{Démarrage à froid du composant d'enrichissement plus long que prévu.} La
        durée de démarrage à froid mesurée sur sept jours de supervision s'est révélée très
        supérieure aux valeurs de référence initialement documentées, avec une moyenne
        d'une trentaine de secondes et des pics proches de la minute et demie --- provoquant
        des erreurs de service début août. La correction a consisté à élargir la fenêtre de
        tolérance de la sonde de disponibilité.
  \item \textbf{Décision assumée de doubler le coût de la base de données pour obtenir une
        haute disponibilité réelle}, prise après un audit de bout en bout, plutôt que de
        conserver l'objectif initial de \og best effort \fg{}.
  \item \textbf{Incident de dérive lors du chiffrement géré par le client sur un secret} :
        une application d'infrastructure partiellement en échec a laissé l'état local
        enregistrer une configuration jamais réellement appliquée côté plateforme cloud.
        L'écart a été détecté par une comparaison explicite entre l'état déclaré et la
        réalité observée, puis corrigé par une opération de resynchronisation avant de
        poursuivre. La leçon retenue : après toute application partiellement en échec, la
        réalité doit être revérifiée avant de replanifier --- ne jamais supposer que
        l'échec a laissé l'état intact.
\end{itemize}

\begin{attentedoc}
Vérifier chaque chiffre ci-dessus contre \texttt{09\_AUDIT\_E2E\_STAGING\_2026-08-07.md}
avant publication --- ce paragraphe reprend l'information disponible au moment de la
rédaction du sommaire de travail, à confirmer mot pour mot.
\end{attentedoc}
